\section{\texorpdfstring{linux network:
\texttt{net\_device}}{linux network: net\_device}}\label{linux-network-net_device}

\section{\texorpdfstring{linux network: \texttt{net\_device} e
\texttt{sk\_buff}}{linux network: net\_device e sk\_buff}}\label{linux-network-net_device-e-sk_buff}

\textbf{fonte}: Linux Kernel Networking - Implementation and Theory.pdf

\subsection{network device drivers}\label{network-device-drivers}

\textbf{Network device drivers}: operano sempre a livello 2 nello stack
di rete, ed e\textquotesingle{} rappresentato nel kernel dalla struttura
\passthrough{\lstinline!net\_device!}. Questa struttura dati contiene
varie informazioni:

\begin{itemize}
\tightlist
\item
  \passthrough{\lstinline!IRQ!}: numero identificativo del dispositivo
\item
  \passthrough{\lstinline!MTU!}: la MTU del dispositivo
\item
  \passthrough{\lstinline!MAC!}: MAC address del dispositivo
\item
  \textbf{nome}: \passthrough{\lstinline!eth0!},
  \passthrough{\lstinline!enp0s3!}, \passthrough{\lstinline!wlan0!},
  ecc...
\item
  \passthrough{\lstinline!net\_device\_ops!}: contiene i puntatori alle
  funzioni per la gestione del dispositivo (start, stop, transmit,
  receive)
\end{itemize}

\textbf{Cosa fa un device driver}?

\begin{itemize}
\tightlist
\item
  \textbf{riceve} i pacchetti destinati al local host, e li inoltra al
  livello di rete, e poi al livello di trasporto.
\item
  \textbf{trasmette} i pacchetti generati dal local host, o fa il
  forward dei pacchetti ricevuti
\end{itemize}

\textbf{routing}: ogni pacchetto in in transito in tutti e due i sensi
triggera un lookup nel sottosistema di routing, che decide cosa bisogna
fare con quel pacchetto. La fase di routing si distingue in 5 fasi,
detti \textbf{hook}, dove possiamo registrare dei \textbf{callback} per
il \textbf{sottosistema netfilter}.

Per esempio l\textquotesingle hook
\passthrough{\lstinline!NF\_INET\_PRE\_ROUTING!}, registra la chiamata
al callback per essere chiamato prima che il lookup del sottosistema d
routing venga effettuato.

\textbf{verdetto}: ogni callback ritorna un verdetto sul pacchetto
decidendo se bisogna accettarlo (\passthrough{\lstinline!NF\_ACCEPT!}),
oppure buttarlo (\passthrough{\lstinline!NF\_DROP!}).

\begin{quote}
con \passthrough{\lstinline!nftables!} applichiamo le regole per il
routing manipolando quindi strutture dati interne al kernel.
\end{quote}

oltre ai callback, \textbf{molti altri fattori} influenzano il modo in
cui il routing viene effettuato. (es: \emph{il TTL viene decrementato e
arriva a 0}).

\subsection{Socket Buffer}\label{socket-buffer}

\begin{lstlisting}[language=C]
struct sk_buff {
    . . .
    struct sock *sk;
    struct net_device *dev;
    . . .
    __u8 pkt_type:3,
    . . .
    __be16 protocol;
    . . .
    sk_buff_data_t tail;
    sk_buff_data_t end;
    unsigned char   *head,
                    *data;
    sk_buff_data_t transport_header;
    sk_buff_data_t network_header;
    sk_buff_data_t mac_header;
    . . .
};
\end{lstlisting}

\begin{itemize}
\tightlist
\item
  \passthrough{\lstinline!head!} e \passthrough{\lstinline!data!}: sono
  puntatori, rispettivamente ad \passthrough{\lstinline!header!} e
  \passthrough{\lstinline!data!} per il layer di rete corrente.
\end{itemize}

Socket buffer e\textquotesingle{} la struttura dati su cui tutti i layer
di rete nel kernel operano.

\subsection{Netfilter}\label{netfilter}

\textbf{fonte}:
\url{https://www.digitalocean.com/community/tutorials/a-deep-dive-into-iptables-and-netfilter-architecture}
\textbf{fonte}:\url{https://linuxhandbook.com/iptables-vs-nftables/}
E\textquotesingle{} un sottosistema del kernel che fornisce le seguenti
funzionalita:

\begin{itemize}
\tightlist
\item
  selezione e filtraggio dei pacchetti
\item
  NAT
\item
  tracciamento delle connessioni
\item
  accounting di rete.
\end{itemize}

Per quanto riguarda il filtraggio dei pacchetti, abbiamo a disposizione
5 hook differenti:

\begin{itemize}
\tightlist
\item
  \passthrough{\lstinline!NF\_INET\_PRE\_ROUTING!}: e\textquotesingle{}
  chiamato nel metodo \passthrough{\lstinline!ip\_rcv()!}, ancora prima
  di entrare nel sottosistema di routing.
\item
  \passthrough{\lstinline!NF\_INET\_LOCAL\_IN!}: e\textquotesingle{}
  chiamato nel momento in cui il pacchetto e\textquotesingle{}
  indirizzato al local host.
\item
  \passthrough{\lstinline!NF\_INET\_FORWARD!}: e\textquotesingle{}
  chiamato dopo che e\textquotesingle{} stato effettuato il lookup nel
  sottosistema di routing
\item
  \passthrough{\lstinline!NF\_INET\_POST\_ROUTING!}: e\textquotesingle{}
  chiamato per i pacchetti che devono essere inoltrati, pochi istanti
  prima che il pacchetto entri fisicamente nel cavo.
\item
  \passthrough{\lstinline!NF\_INET\_LOCAL\_OUT!}: e\textquotesingle{}
  chiamato per i pacchetti creati nel local host, prima che escano.
\end{itemize}

Molte utility, comunemente usate, si basa sul sottosistema netfilter:

\begin{itemize}
\tightlist
\item
  \passthrough{\lstinline!iptables!}: agisce come front-end per
  netfilter ed implementa un \textbf{firewall}. Fornisce
  un\textquotesingle interfaccia per aggiungere, eliminare, regole
  netfilter, aggiungere tabelle, ecc...
\item
  \passthrough{\lstinline!nftables!}: e\textquotesingle{} sviluppato da
  \textbf{netfilter} (l\textquotesingle organizzazione dietro al
  sottosistema netfilter). E\textquotesingle{} un rimpiazzo per iptables
  che soffre di problemi di scalabilita e performance.
\end{itemize}

\textbf{iptables}: utilizza le \textbf{tables} per organizzare le regole
da applicare, suddividendo le regole in classi di \textbf{table}.

\begin{itemize}
\tightlist
\item
  \passthrough{\lstinline!nat!} table: usata per implementare le regole
  NAT, ossia si modificano sorgente e destinazione dei pacchetti in
  transito nello stack di rete per modificarne il comportamento durante
  il routing.
\item
  \passthrough{\lstinline!mangle!} table: altera gli header IP. (es:
  modificare il TTL, marcare il pacchetto)
\item
  \passthrough{\lstinline!raw!} table.
\item
  \passthrough{\lstinline!security!} table: per inserire contesto
  SELinux all\textquotesingle interno dei pacchetti.
\end{itemize}

\textbf{nftables}: svolege la stessa funzione di iptables, ma usando una
sintassi differente. rispetto ad iptables, nftables e\textquotesingle{}
piu efficiente in memoria ed in CPU ed e\textquotesingle{} piu
flessibile.

\subsection{Connection Tracking}\label{connection-tracking}

Spesso le connessioni a livello applicativo hanno una connessione
mantenuta nel tempo ed una serie di eventi che determinano la storia
della comunicazione. Connection Tracking permette al kernel di tenere
traccia delle sessioni.

Una Connection Tracking entry nel kernel appare cosi:

\begin{lstlisting}[language=C]
struct nf_conn {
    /* Usage count in here is 1 for hash table/destruct timer, 1 per skb,
    plus 1 for any connection(s) we are `master' for */
    struct nf_conntrack ct_general;
    spinlock_t lock;
    /* XXX should I move this to the tail ? - Y.K */
    /* These are my tuples; original and reply */
    struct nf_conntrack_tuple_hash tuplehash[IP_CT_DIR_MAX];
    /* Have we seen traffic both ways yet? (bitset) */
    unsigned long status;
    /* If we were expected by an expectation, this will be it */
    struct nf_conn *master;
    /* Timer function; drops refcnt when it goes off. */
    struct timer_list timeout;
    . . .
    /* Extensions */
    struct nf_ct_ext *ext;
#ifdef CONFIG_NET_NS
    struct net *ct_net;
#endif
    /* Storage reserved for other modules, must be the last member */
    union nf_conntrack_proto proto;
};
\end{lstlisting}

e\textquotesingle{} la struttura dati fondamentale per implementare il
NAT. Dunque per ogni connessione che deve usufruire del NAT deve tenere
in memoria una connection Tracking entry.

\section{TODO: NAT e masquarade}\label{todo-nat-e-masquarade}

\section{bridge}\label{bridge}

\textbf{fonte}:
\href{https://books.google.it/books?hl=en&lr=&id=ALapr7CvAKkC&oi=fnd&pg=PR5&dq=linux+bridges+L2&ots=gYfwjUhYJn&sig=9yrrq3kOgtMhkfkJiZP9hF1BgRw&redir_esc=y\#v=onepage&q&f=true}{libro!}

Un bridge e\textquotesingle{} un dispositivo che capisce il protocollo
del livello di collegamento, effettua l\textquotesingle inoltro dei dati
copiandoli frame per frame e deve essere in grado di mantenere nel suo
buffer almeno un frame per volta.

Un bridge assegna un\textquotesingle indirizzo a livello di collegamento
ad ognuna delle sue interfacce ed effettua il forward verso di queste se
il bridge non e\textquotesingle{} il destinatario del pacchetto.

\begin{itemize}
\tightlist
\item
  il bridge, se implementa delle funzionalita al livello L3 potrebbe
  passare il frame al livello superiore e processarlo.
\item
  oppure potrebbe passare il frame ad un protocol handler
\end{itemize}

\textbf{utilizzo}: un bridge viene di solito usato per unire due LAN
differenti. !{[}{]}(Pasted image 20261001233828.png)

\textbf{definizione di host}: un qualsiasi device che opera e capisce ad
un livello più alto del layer di collegamento.

\textbf{address learning}: un bridge fa \emph{passive learning}.
Ipotizziamo che \passthrough{\lstinline!A!} invii un frame verso
\passthrough{\lstinline!B!} e che \passthrough{\lstinline!B!} lo riceva
poiche \passthrough{\lstinline!A!} e \passthrough{\lstinline!B!} si
trovano nella stessa LAN1. Il bridge, riceve anche lui una copia del
frame, sulla porta LAN1 e lo inoltra verso LAN2 (ance se non serviva)

\begin{itemize}
\tightlist
\item
  \passthrough{\lstinline!learning!}: ora il bridge sa che sulla porta
  LAN1 si trova \passthrough{\lstinline!A!}
\item
  ora se \passthrough{\lstinline!B!} invia un frame verso
  \passthrough{\lstinline!A!}, il bridge sa che
  \passthrough{\lstinline!B!} si trova nella stessa rete di
  \passthrough{\lstinline!A!} e non deve copiare il frame su LAN2.
\end{itemize}

\textbf{broadcast}: se un bridge riceve un frame con destinazione
FF:FF:FF:FF:FF:FF, il frame viene copiato su tutte le porte.

\textbf{aging}: quando imparo l\textquotesingle indirizzo di un nuovo
host, aggiungo un timer di aging. Quando il timer scade, rimuovo
l\textquotesingle indirizzo dalla tabella, se arriva un nuovo pacchetto
prima della scadenza, estendo il timer.

\section{linux bridge}\label{linux-bridge}

E\textquotesingle{} un device virtuale: finché non viene collegato ad
altri device non puo\textquotesingle{} ne ricevere e ne trasmettere
frame.

Supponiamo di voler implementare questa topologia, nel kernel:
!{[}{]}(Pasted image 20261002000350.png)

Ossia, data una macchina con le tre interfacce fisiche
\passthrough{\lstinline!eth0, eth1, eth2!}vogliamo aggregare le prime
due in un\textquotesingle unica interfaccia in modo che questa possa
instradare verso \passthrough{\lstinline!eth2!}.

Il kernel dovra:

\begin{itemize}
\tightlist
\item
  al livello del router, vedere le due sottoreti
\item
  effettuare il bridge tra eth0 ed eth1 anche se ho piu interfacce in
  pratica
\end{itemize}

\textbf{enslaving}: quando si crea un bridge, bisogna assegnare a questo
delle interfacce dette \textbf{enslaved}.

Dunque possiamo creare il bridge \passthrough{\lstinline!br0!} a cui
assegnamo le interfacce \passthrough{\lstinline!eth0!} ed
\passthrough{\lstinline!eth1!}. Dato che \passthrough{\lstinline!eth0!}
ed \passthrough{\lstinline!eth1!} sono interfacce al bridge, non hanno
bisogno di essere viste a livello 3, e dunque non hanno
un\textquotesingle indirizzo IP.

Si assegna un\textquotesingle indirizzo IP all bridge per renderlo
partecipe nella sotto rete, in modo che possa ricevere i pachetti da cui
estrarre i frame da inoltrare.

In pratica, quando tento di trasmettere su
\passthrough{\lstinline!br0!}, il device driver di
\passthrough{\lstinline!br0!} gestisce i device \textbf{enslaved}
effettuando il forwarding dei frame.

Posso eventualmente nell\textquotesingle host, assegnare
un\textquotesingle indirizzo IP anche alle interfacce enslaved. Da quel
momento il kernel linux si occuperà di instradare verso
l\textquotesingle interfaccia si attraverso
\passthrough{\lstinline!br0!}, che attraverso il layer 3.

\begin{quote}
{[}!error{]} eventualmente approfondire Capitolo 8, 10 e 11 di
Understanding Linux Network Internals.pdf
\end{quote}

\subsection{STP}\label{stp}

Lo Spanning Tree Protocol serve a prevenire loop
all\textquotesingle interno della rete nel kernel. Con STP abilitato, i
bridge si scambiano Bridge Protocol Data Units in modo da identificare
un bridge root e l\textquotesingle interfaccia da bloccare per rendere
la topologia della rete libera da cicli.

\textbf{cicli}: all\textquotesingle interno di una rete, a livello L2
puo\textquotesingle{} causare ondate di messaggi.

\section{cosa approfondire}\label{cosa-approfondire}

\subsection{shadow bridge}\label{shadow-bridge}

\begin{itemize}
\tightlist
\item
  differenze tra switch L2 e router L3
\item
  risoulzione next-hop e forwarding information base: Quando un
  container vuole comunicare con un IP esterno alla propria subnet (es.
  DNS \passthrough{\lstinline!1.1.1.1!} o repository di pacchetti), la
  Forwarding Information Base (FIB) del container seleziona il default
  gateway (\passthrough{\lstinline!10.k.1.1!}) e invia una richiesta ARP
  broadcast (\passthrough{\lstinline!who-has 10.k.1.1!}).
\item
  FIB collision: accade quando sistenta di assegnare
  \passthrough{\lstinline!10.k.1.1!} a due bridge
\end{itemize}

\subsection{dual nic}\label{dual-nic}

\begin{itemize}
\item
  multi-homing: due interfacce che instradando verso gateway distinti
\item
  \begin{enumerate}
  \def\labelenumi{\arabic{enumi}.}
  \setcounter{enumi}{2}
  \tightlist
  \item
    \textbf{Policy Routing \& Source-Based Routing
    (\passthrough{\lstinline!ip rule!})}:
  \end{enumerate}

  \begin{itemize}
  \tightlist
  \item
    Per far funzionare due interfacce con due default route,
    bisognerebbe iniettare regole di routing avanzato nel container
    basate sull\textquotesingle IP sorgente
    (\passthrough{\lstinline!ip rule add from ... table ...!}). Questo
    viola la semplicità di configurazione e sporca
    l\textquotesingle ambiente del concorrente.
  \item
    fedalta topologica: aggiungere una seconda interfaccia sporca il
    comportamento di demoni che usano l\textquotesingle interfaccia
    \passthrough{\lstinline!0.0.0.0!}
  \item
    \passthrough{\lstinline!veth!}: numero raddoppiato.
  \end{itemize}
\end{itemize}

\subsection{PID namespace separato}\label{pid-namespace-separato}

\begin{enumerate}
\def\labelenumi{\arabic{enumi}.}
\tightlist
\item
  PID vs mount namespace: isolare il PID namespace non isola il
  filesystem \passthrough{\lstinline!/proc!}
\item
  \passthrough{\lstinline!pspy!} usa \passthrough{\lstinline!inotify!}:
  qualsiasi directory creata in \passthrough{\lstinline!/proc!} viene
  cattuarata da \passthrough{\lstinline!pspy!}
\item
  file system unico: il partecipante puo analizzare i log ed estrarre
  ventualmente il payload
\item
  Network Sniffing su loopback: se i processi condividono lo stesso
  network namespace, allora il partecipante puo\textquotesingle{}
  analizzare il payload se transita su un\textquotesingle interfaccia di
  rete.
\item
  il payload potrebbe far carashare il server in esecuzione
  nell\textquotesingle ambiente dello sfidante
\end{enumerate}

\subsection{network namespace
separato}\label{network-namespace-separato}

\begin{enumerate}
\def\labelenumi{\arabic{enumi}.}
\tightlist
\item
  \passthrough{\lstinline!CLONE\_NEWNET!},
  \passthrough{\lstinline!ip netns!} crea un\textquotesingle istanza
  isolata dello stack di rete. permette a due brdige di avere lo stesso
  IP
\item
  \passthrough{\lstinline!veth!}: per dare connessione serve una coppia
  di interfacce \passthrough{\lstinline!veth!} con
  \passthrough{\lstinline!ip forwarding!} e
  \passthrough{\lstinline!SNAT/MASQUERADE!}
\item
  creazione di interfacce non tracciate da incus, disallinamento
\item
  \passthrough{\lstinline!rtnl\_lock!}: creare e distruggere network
  namespace comporta il lock sul routing globale
\end{enumerate}

\subsection{esecuzione in memoryu}\label{esecuzione-in-memoryu}

\begin{itemize}
\tightlist
\item
  \textbf{dynamic linker hooking}: sovrascrivere la libreria
  \passthrough{\lstinline!C!} per filtrare e nascondere
  \passthrough{\lstinline!/proc!}

  \begin{itemize}
  \tightlist
  \item
    funziona solo con binari linkati dinamicamente
  \item
    \passthrough{\lstinline!/etc/ld.so.prelaod!} puo\textquotesingle{}
    essere modificato
  \end{itemize}
\item
  \textbf{fileless execution}: \passthrough{\lstinline!memfd\_create!},
  crea un descrittore di file anonimo residente unicamente in RAM
  (memoria volatile protetta).

  \begin{itemize}
  \tightlist
  \item
    \passthrough{\lstinline!fexecve!} lo esegue
  \end{itemize}
\item
  usando filtri \passthrough{\lstinline!eBPF!} posso a livello kernel
  intercettare le chiamate di sistema
\item
  principio di saltxzer e schroeder:

  \begin{itemize}
  \tightlist
  \item
    l\textquotesingle isolamento deve basarsi su barriere architetturali
    matematicamente verificate (hypervisor, namespaces, CoW snapshots),
    non su trucchi di evasione ("security through obscurity").

    \begin{itemize}
    \tightlist
    \item
      L\textquotesingle occultamento non previene il salvataggio dei log
      applicativi, non isola il traffico di rete e non evita il crash
      distruttivo del servizio.
    \end{itemize}
  \end{itemize}
\end{itemize}
